% =========================================================================
% ปฏิบัติการที่ 2: วิทัศน์คอมพิวเตอร์และการตรวจจับโครงร่างมือ
% =========================================================================

\chapter{วิทัศน์คอมพิวเตอร์และการตรวจจับโครงร่างมือ}
\label{ch:lab02}

\noindent{\large\bfseries\color{cyberblue} การประมวลผลวิดีโอสตรีมสดและสกัดพิกัดมือ 21 จุดร่วมด้วย MediaPipe Hands}

\vspace{0.4cm}

\begin{labobjectivebox}{วิทัศน์คอมพิวเตอร์และการตรวจจับโครงร่างมือ}
\textbf{วัตถุประสงค์การเรียนรู้}
\begin{enumerate}[leftmargin=1.8cm, itemsep=2pt]
    \item เข้าใจทฤษฎีโครงข่ายประสาทเทียมตรวจจับจุดสำคัญของมือ (Hand Landmark Topology)
    \item สามารถเชื่อมต่อสตรีมวิดีโอสดจากเว็บแคมเข้าสู่ไปป์ไลน์ MediaPipe Hands
    \item สามารถสกัดและวาดจุดพิกัด 21 ตำแหน่งลงบนผืนผ้าใบแบบเรียลไทม์
\end{enumerate}

\vspace{4pt}
\textbf{เครื่องมือและอุปกรณ์จำเป็น} \enspace เว็บแคมความละเอียด 720p ขึ้นไป, MediaPipe Hands API (@mediapipe/camera\_utils, @mediapipe/hands)
\end{labobjectivebox}

\section{หลักการและทฤษฎีพื้นฐาน}

เทคโนโลยีการตรวจจับโครงกระดูกมือ (Hand Pose Estimation) ของ MediaPipe อาศัยแบบจำลองการเรียนรู้เชิงลึกแบบสองขั้นตอน (Two-stage Deep Learning Pipeline) ได้แก่:
1. โมเดลตรวจจับฝ่ามือ (Palm Detector Model): ทำหน้าที่ค้นหาขอบเขตของฝ่ามือจากภาพวิดีโอทั้งภาพ (Bounding Box Detection) ซึ่งเป็นขั้นตอนที่ตรวจจับได้ง่ายเนื่องจากฝ่ามือมีลักษณะสัณฐานคงที่และบดบังตัวเองน้อยกว่านิ้วมือ
2. โมเดลจุดสังเกตสำคัญ (Hand Landmark Model): เมื่อทราบตำแหน่งฝ่ามือ โมเดลจะตัดภาพเฉพาะบริเวณมือมาประมวลผล เพื่อพยากรณ์พิกัด 3 มิติ $(x, y, z)$ ของจุดร่วมทั้ง 21 จุดสำคัญ (Landmarks) บนมืออย่างแม่นยำ

พิกัด $x$ และ $y$ จะถูกปรับให้เป็นสเกลมาตรฐานสัมพัทธ์ $[0.0, 1.0]$ ตามความกว้างและความสูงของภาพ ส่วนพิกัด $z$ แสดงความลึกสัมพัทธ์ (Relative Depth) โดยมีจุดอ้างอิงอยู่ที่ข้อมือ (Landmark 0: Wrist) ซึ่งจุดสำคัญเหล่านี้ถูกจัดกลุ่มเป็น 5 นิ้ว ได้แก่ นิ้วหัวแม่มือ (Landmarks 1-4), นิ้วชี้ (Landmarks 5-8), นิ้วกลาง (Landmarks 9-12), นิ้วนาง (Landmarks 13-16) และนิ้วก้อย (Landmarks 17-20)

\begin{conceptbox}{สาระสำคัญของปฏิบัติการที่ 2}
การเข้าใจโครงสร้างทางคณิตศาสตร์และการประมวลผลเชิงพื้นที่ ช่วยให้นักศึกษาสามารถออกแบบระบบเสมือนจริงที่ตอบสนองต่อผู้ใช้งานได้อย่างเป็นธรรมชาติและแม่นยำสูง ทั้งยังสามารถต่อยอดสู่การพัฒนาแอปพลิเคชันทางวิทยาศาสตร์และอุตสาหกรรมได้อย่างมีประสิทธิภาพ
\end{conceptbox}

\section{ขั้นตอนการปฏิบัติการและการทดลอง}

ให้นักศึกษาดำเนินการสร้างไฟล์โครงการและเขียนคำสั่งตามลำดับขั้นตอนต่อไปนี้ โดยตรวจสอบความถูกต้องของไลบรารีที่เรียกใช้งานและเส้นทางไฟล์

\begin{codebox}{โครงสร้างโค้ดหลักของปฏิบัติการที่ 2}
\begin{lstlisting}[language=HTML]
// กำหนดค่าและเริ่มต้นใช้งาน MediaPipe Hands
const videoElement = document.getElementById('webcam-video');
const canvasElement = document.getElementById('output-canvas');
const canvasCtx = canvasElement.getContext('2d');

const hands = new Hands({
    locateFile: (file) => `https://cdn.jsdelivr.net/npm/@mediapipe/hands/${file}`
});

hands.setOptions({
    maxNumHands: 2,
    modelComplexity: 1,
    minDetectionConfidence: 0.7,
    minTrackingConfidence: 0.7
});

hands.onResults((results) => {
    canvasCtx.save();
    canvasCtx.clearRect(0, 0, canvasElement.width, canvasElement.height);
    canvasCtx.drawImage(results.image, 0, 0, canvasElement.width, canvasElement.height);

    if (results.multiHandLandmarks) {
        for (const landmarks of results.multiHandLandmarks) {
            // วาดเส้นเชื่อมต่อกระดูกและจุดพิกัด 21 ตำแหน่ง
            drawConnectors(canvasCtx, landmarks, HAND_CONNECTIONS, {color: '#06b6d4', lineWidth: 3});
            drawLandmarks(canvasCtx, landmarks, {color: '#f59e0b', lineWidth: 1, radius: 4});
            
            // สกัดพิกัดปลายนิ้วชี้ (Landmark 8) และนิ้วโป้ง (Landmark 4)
            const indexTip = landmarks[8];
            const thumbTip = landmarks[4];
            console.log(`Index: (${indexTip.x.toFixed(3)}, ${indexTip.y.toFixed(3)}, ${indexTip.z.toFixed(3)})`);
        }
    }
    canvasCtx.restore();
});

const camera = new Camera(videoElement, {
    onFrame: async () => { await hands.send({image: videoElement}); },
    width: 1280, height: 720
});
camera.start();
\end{lstlisting}
\end{codebox}

\section{การทดสอบระบบและการบันทึกผล}

เมื่อรันโปรแกรมบนเซิร์ฟเวอร์จำลอง (Local Development Server) ให้ทำการทดสอบการทำงานตามเกณฑ์ประเมินในตารางต่อไปนี้ พร้อมบันทึกผลการสังเกตการณ์ลงในรายงานผลการทดลอง

\begin{table}[htbp]
\centering
\small
\begin{tabularx}{\textwidth}{c X c Y}
\toprule
\textbf{ลำดับ} & \textbf{รายการทดสอบและพฤติกรรมที่คาดหวัง} & \textbf{เกณฑ์ผ่าน} & \textbf{ผลการทดสอบ} \\
\midrule
1 & การเริ่มต้นระบบและการเข้าถึงสิทธิ์ฮาร์ดแวร์ & ภายใน 3 วินาที & ผ่าน / ไม่ผ่าน \\
2 & ความเสถียรของการตรวจจับและการคำนวณพิกัด & อัตราความแม่นยำ $> 90\%$ & ผ่าน / ไม่ผ่าน \\
3 & อัตราการแสดงผลกราฟิกต่อเนื่อง (Framerate) & ไม่ต่ำกว่า 55 FPS & ผ่าน / ไม่ผ่าน \\
4 & การตอบสนองต่อปฏิสัมพันธ์เชิงพื้นที่ & ความหน่วง $< 40$ ms & ผ่าน / ไม่ผ่าน \\
\bottomrule
\end{tabularx}
\caption{ตารางประเมินผลการทำงานของระบบในปฏิบัติการที่ 2}
\label{tab:eval_lab02}
\end{table}

\section{ภารกิจท้าทายและการต่อยอด}

\begin{activitybox}{กิจกรรมส่งเสริมทักษะขั้นสูง}
ให้นักศึกษาเขียนโปรแกรมเพิ่มเติมเพื่อจำแนกมือซ้ายและมือขวาจากข้อมูล `results.multiHandedness` และแสดงข้อความระบุชื่อมือพร้อมค่าความเชื่อมั่น (Confidence Score) บนผืนผ้าใบ
\end{activitybox}

\section{คำถามท้ายการทดลอง}

ให้นักศึกษาตอบคำถามเชิงวิเคราะห์ต่อไปนี้ลงในสมุดบันทึกปฏิบัติการ

\begin{enumerate}[leftmargin=1.8cm, itemsep=6pt]
    \item เหตุใดสถาปัตยกรรมของ MediaPipe จึงเลือกใช้การตรวจจับฝ่ามือก่อนการพยากรณ์จุดสำคัญ 21 จุดแทนการค้นหานิ้วมือโดยตรง
    \item พิกัดแกน $z$ ที่รายงานโดยโมเดลจุดสังเกตสำคัญของ MediaPipe มีความหมายทางกายภาพอย่างไรและอ้างอิงกับจุดใด
    \item หากสภาพแสงในห้องทดลองมีความสว่างน้อย จะส่งผลต่อค่า minDetectionConfidence และอัตราการกระตุกของเฟรมอย่างไร
\end{enumerate}

